% !PRINTED 586
\weberpage{586}

%<<CONT
$\equiv 1$ (mod 4), also $a^2+b^2\equiv 1$ (mod 4) sein muss. Der Satz 3.,
der übrigens den zweiten in sich schliesst, wird so bewiesen.

Angenommen, es sei $a^2+b^2$, aber nicht $b$, durch $q$ theilbar,
so bestimmen wir $b'$ aus der Congruenz $bb'\equiv 1$ (mod $q$) und
erhalten aus $a^2+b^2\equiv 0$ (mod $q$)
\[
(ab')^2\equiv -1\ (\mathrm{mod}\ q).
\]
Dies ist aber unmöglich, da nach \S.\,138, 4. für jede Primzahl $q$
die Zahl $-1$ quadratischer Nichtrest ist.

Der Körper $R(i)$ ist der Inbegriff aller Zahlen von der
Form $x+yi$, wenn $x$, $y$ ganze oder gebrochene rationale Zahlen
sind.

\gesperrt{Die Zahlen} $a+bi$, \gesperrt{in denen} $a$ \gesperrt{und} $b$ \gesperrt{ganze Zahlen
in} $R$ \gesperrt{sind, heissen die ganzen Zahlen des Körpers} $R(i)$.

Das Product zweier conjugirter Zahlen
\[
\xi=x+yi,\quad \xi'=x-yi,
\]
also
\[
\xi\xi'=x^2+y^2
\]
heisst die Norm von $\xi$ und wird mit $N(\xi)$ bezeichnet. Die Norm
einer nicht verschwindenden Zahl $\xi$ ist eine positive rationale
Zahl und die Norm einer ganzen Zahl ist eine ganze Zahl. Die
Norm eines Productes oder eines Quotienten ist gleich dem Product
 oder dem Quotienten der Normen.

\gesperrt{Eine ganze Zahl, deren Norm gleich 1 ist, heisst
eine Einheit.}

Da $a^2+b^2$ (für ganzzahlige $a$, $b$) nur dann $=1$ sein kann,
wenn $a=\pm 1$, $b=0$ oder $a=0$, $b=\pm 1$ ist, so giebt es in
$R(i)$ nur die vier Einheiten
\[
+1,\quad -1,\quad +i,\quad -i.
\]
Der reciproke Werth einer Einheit ist auch eine Einheit.

Zwei Zahlen, von denen die eine aus der anderen durch
Multiplication mit einer Einheit entsteht, heissen \gesperrt{associirte
Zahlen}.

Jede complexe Zahl gehört zu einem System von vier associirten
 Zahlen
\[
a+bi,\quad -a-bi,\quad -b+ai,\quad b-ai.
\]
Summe, Differenz und Product zweier ganzer Zahlen sind
wieder ganze Zahlen.

\gesperrt{Eine ganze Zahl} $\alpha$ \gesperrt{heisst durch eine ganze Zahl} $\beta$
\gesperrt{theilbar, wenn eine dritte ganze Zahl} $\gamma$ \gesperrt{existirt, so
dass} $\alpha=\beta\gamma$ \gesperrt{ist.}
